\section{Discussion and limitations}

\subsection{What the measurement-to-control gap means}
Across two EEG representations, conditional domain leakage is \textbf{measurable} by the score-Fisher
diagnostics but is \textbf{not a sufficient control target} for cross-subject generalization. The way control
fails is representation-dependent: in a high-dimensional SPD/LogEig latent the leakage is redundantly
re-encoded and is not low-rank removable, and a global penalty removes it only by collapsing the
representation; in a compact convolutional latent the leakage \emph{is} low-rank removable (linearly) on the
source side, yet a separate collapse-free global penalty that reduces it does not improve target accuracy
(the frozen selective erasers' target effect is left to a deployment test). The common thread is a gap between \textbf{localizing/measuring}
leakage and \textbf{controlling} it: neither localization nor removal is, by itself, evidence of a
generalization benefit. This reframes selective conditional invariance as a \textbf{certified intervention with
refusal} --- delete only when a source-risk certificate permits, and abstain otherwise.

A secondary, mechanistic observation concerns \emph{how} the global penalty fails at large $\lambda$. The collapse is an
\textbf{objective-scaling / training-basin pathology} --- the latent contracts to the origin, which trivially
satisfies the penalty --- not a gradient explosion and not a smooth geometric over-compression (the apparently
stable effective rank is a scale-invariant-metric artifact). It is preventable by warm-up scheduling, but
preventing it returns the leakage to baseline, so the apparent de-domaining was an artifact of the collapse.

\subsection{Why leakage removal need not improve generalization}
The EEGNet result is the cleanest case: leakage is genuinely (if partially) removable at no task cost, yet
mean LOSO target accuracy is flat-to-worse and uncorrelated with the amount of leakage removed. Two
readings are consistent with our evidence. First, conditional domain leakage $I(Z;D|Y)$ may be \textbf{only weakly
coupled to the component of cross-subject error that matters} on this paradigm --- subject identity is
\emph{present and decodable} without being the \emph{cause} of the transfer gap. Second, removing a linearly-decodable
slice of identity can leave a \textbf{nonlinear residual} (we observe one) and can also remove incidentally
task-correlated variance, so the net effect on the target task is neutral-to-negative. In either reading,
leakage decodability is a property of the representation, not a certificate of a causal generalization
lever --- which is exactly why we do not equate removal with improvement.

\subsection{Limitations}
\begin{itemize}
\item \textbf{Single dataset.} All EEG results are on BCI-IV-2a (BNCI2014\_001); we do not claim generality to other
  paradigms, montages, or session structures.
\item \textbf{Two backbones only.} TSMNet and EEGNet; the cross-representation contrast rests on two points.
\item \textbf{Dimension-vs-type: largely capacity, with a residual architecture effect at high $d_z$.} The original
  two-point comparison confounded architecture with latent dimension (SPD/$210$ vs conv/$16$). The
  architecture$\times$dimension factorial (\S4.5, Fig.~\ref{fig:capacity}; 3 seeds, fold-cluster CIs) shows the
  nonlinear erasure residual is primarily a function of latent \emph{dimension} --- it rises with $d_z$ in both
  architectures, the two coincide at low $d_z$, and matching dimension removes about two-thirds of the raw gap ---
  but \emph{not} purely: at matched $d_z{=}210$ the SPD/LogEig latent retains significantly more recoverable
  subject identity than the convolutional one ($+0.11$, CI $[0.09,0.13]$). So the confound is largely, not fully,
  explained by capacity. Remaining caveats: a single dataset (BCI-IV-2a) and the LDA cap from $8$ source subjects.
\item \textbf{No end-to-end TOS training.} The EEG study is a frozen-feature pilot: it measures whether \emph{removing}
  localized leakage from a fixed representation helps; it does not show that end-to-end training with a
  certified penalty cannot help.
\item \textbf{Certified deletion mostly abstains.} On synthetic controls the power-certified gate is conservative
  (frequent abstention); on EEG it abstains on TSMNet and is only diagnostic on EEGNet. By design --- we
  report a certification framework with refusal, not a deletion method that is on by default.
\item \textbf{Leakage removal is not shown to be a causal DG mechanism.} Even where removal succeeds (EEGNet), target
  accuracy does not improve; we make no causal claim linking conditional domain leakage to the transfer gap.
\item \emph{(Secondary)} Probe/estimator dependence: the nonlinear residual is reported under an MLP probe; the
  per-epoch encoder gradient is a between-epoch diagnostic proxy; the Fisher/LDA analysis is capped at 7
  directions by 8 source subjects under LOSO.
\end{itemize}

\subsection{Implications for EEG invariance methods}
Three practical implications follow. (i) \textbf{Report the mechanism, not just the metric:} a global
conditional-invariance penalty that ``removes leakage'' should be checked for representation collapse (feature
norm, not effective rank), because a falling leakage proxy can coincide with a degenerate representation.
(ii) \textbf{Treat selective deletion as a certified intervention with refusal:} localize, test the conditional
task risk, and default to the identity map when a safe, useful deletion cannot be certified at the available
sample size. (iii) \textbf{Do not treat leakage decodability as a generalization target:} measure whether
removing it actually moves target accuracy before claiming a DG benefit. The natural next steps ---
multi-dataset hardening of the (largely) capacity-mediated verdict (the multi-seed factorial is complete on 2a), source-OOD benefit gating on a large-$N$ dataset, and folding the certificate into end-to-end training with
refusal --- are deferred to future work.
